\chapter{Conclusion and Future Work}
\label{ch:conclusion}

This chapter provides a summary of the development of \textbf{JurisChain}, a LangChain-powered RAG framework designed for Indian Family Law. It reflects on the technical contributions, identifies the operational boundaries encountered during development, and outlines potential trajectories for further research.

\section{Summary of Contributions}
The primary objective of this research was to address the inefficiencies within the Indian judicial process, specifically concerning the backlog of legal documents and the difficulty of manual precedent retrieval. This work successfully delivered:

\begin{itemize}
	\item \textbf{Dataset Creation:} Created a Parquet and Vector Dataset for use in training any Legal models in Indian jurisprudence.
	\item \textbf{System Framework:} The development of \textit{JurisChain}, an end-to-end modular pipeline that integrates LangChain for workflow orchestration and Chroma DB for vector-based semantic retrieval.
	\item \textbf{API-Driven Inference:} A seamless integration with the Gemini API to handle complex reasoning tasks, including evidence scrutinization, verdict prediction, and document summarization, effectively bypassing the token-limitations of local encoders.
	\item \textbf{Legal Synthesis:} A functional system capable of mapping IPC/BNS sections directly from raw case documents, providing a transparent and evidence-based layer for judicial decision support.
	\item \textbf{Academic Validation:} A comprehensive methodology that demonstrates how retrieval-augmented generation (RAG) can mitigate the challenges of long-document processing in the Indian legal landscape, offering a reproducible model for future academic and professional legal tools.
\end{itemize}

\section{Limitations of the Current System}
Despite the successful implementation and performance metrics of the system, several constraints were identified:

\begin{itemize}
	\item \textbf{Dependency on External APIs:} While the reliance on the Gemini API provided superior reasoning and context windows, it introduces operational dependency on external service availability and potential latency issues during high-load periods.
	\item \textbf{Indic Language Translation Bottlenecks:} The current translation module exhibits varying performance with regional vernacular inputs (Marathi/Hindi), where domain-specific legal jargon often suffers from semantic degradation during translation.
	\item \textbf{Statutory Complexity:} The system occasionally struggles with nested, cross-statute references within highly complex or convoluted legal complaints, resulting in suboptimal IPC/BNS mapping accuracy in rare instances.
	\item \textbf{OCR and Document Noise:} The system's performance remains inherently tied to the quality of the source documents. Scanned PDFs with low-resolution text or poor layout formatting can impede the initial extraction and synthesis stages.
	\item \textbf{Test Sample Procurement:} Procuring real world samples containing evidences, petitions, responses, etc from both sides proves to be extremely difficult due to the highly sensitive nature of Family court cases.
	\item \textbf{High Complexity and Relevance of Emotions:} Family court cases can be extremely complex sometimes containing 100s of different parties with their own reasoning, actions and circumstances. LLMs also do not consider the emotions of the entangled parties leading to loss of nuance/context and poor decisions.
\end{itemize}

\section{Future Research Directions}
To build upon the current research, the following avenues are proposed:

\begin{itemize}
	\item \textbf{Domain-Specific Translation Fine-Tuning:} Future iterations could involve fine-tuning open-source Indic language models specifically on parallel legal corpora to enhance the precision of Marathi and Hindi translation, reducing jargon loss.
	\item \textbf{Agentic Pipeline Expansion:} Integrating more robust "Agentic" workflows using LangGraph to allow the system to perform multi-step, iterative reasoning, enabling it to cross-question its own predictions against multiple retrieved precedents.
	\item \textbf{Domain Transfer:} Transferring the model/Pipeline/Vector DB from Indian Family courts to Criminal, Civil, Industrial courts, etc.
	\item \textbf{Larger Models/Quantum Models:} Training LLMs with over 100B parameters or working with Quantum algorithms/models to better handle Family court complexity.
	\item \textbf{Multi-modal Evidence Support:} Expanding the scope of \textit{JurisChain} to ingest multi-modal evidence, such as audio transcripts from courtroom hearings or video-based evidence, to provide a more holistic view of case facts.
\end{itemize}